% =============================================================================
% tema/boxes.tex - caixas semanticas, ferramentas de tabela e macros de autor
% -----------------------------------------------------------------------------
% Decisao herdada da serie: as caixas de conteudo sao construidas sobre o block
% nativo do Beamer, nao sobre tcolorbox, para que raio de canto, respiro interno
% e tipografia de titulo sejam os mesmos em toda a serie. O arredondamento e a
% sombra estao no preambulo, com \setbeamertemplate{blocks}[rounded][shadow],
% e nao aqui.
%
% \CRnewblock tem QUATRO argumentos. O marcador de itemizacao herda a cor pela
% definicao de "local structure" e "item" com a cor do TITULO da caixa.
% =============================================================================

% #1 nome do ambiente | #2 cor do titulo | #3 cor do fundo | #4 cor do texto do titulo
\newcommand{\CRnewblock}[4]{%
  \newenvironment{#1}[1]{%
    \begingroup
    \setbeamercolor{block title}{bg=#2, fg=#4}%
    \setbeamercolor{block body}{bg=#3, fg=black}%
    \setbeamercolor{local structure}{fg=#2}%
    \setbeamercolor{item}{fg=#2}%
    \begin{block}{##1}%
  }{%
    \end{block}%
    \endgroup
  }%
}

% =============================================================================
% AS SEIS CAIXAS. Uma por funcao do mapa de cor do briefing. Nenhuma outra cor.
%   azul     avaliacao e entrevista
%   ambar    mudanca recente e janelas de vigilancia
%   bordo    acao imediata, disposicao, remocao, ambiente protegido
%   verde    suporte, protecao, plano de seguranca, continuidade
%   grafite  normativa, documento e registro
%   roxo     farmacologia
% =============================================================================
\CRnewblock{CRbaseavaliar}{campoazul!88}{campoazul!8}{white}
\CRnewblock{CRbasemudanca}{campoambar}{campoambar!10}{white}
\CRnewblock{CRbaseagir}{campobordo}{campobordo!7}{white}
\CRnewblock{CRbaseproteger}{campoverde}{campoverde!8}{white}
\CRnewblock{CRbasenorma}{campografite}{campografite!7}{white}
\CRnewblock{CRbasefarmaco}{camporoxo}{camporoxo!8}{white}

% =============================================================================
% VOCABULARIO DE ICONE
% -----------------------------------------------------------------------------
% Um icone por FUNCAO, sempre o mesmo, nunca outro. O icone nao decora: ele diz,
% antes da leitura, a que familia aquela caixa pertence, do mesmo modo que a cor.
% Cor e icone andam juntos e nenhum dos dois aparece sozinho.
%
% Todos os nomes usados aqui constam da Comprehensive LaTeX Symbol List e foram
% conferidos um a um contra a instalacao (ferramentas/simbolos.py).
% =============================================================================
\newcommand{\iAvaliar}{\faComments}          % avaliacao e entrevista
\newcommand{\iMudanca}{\faClock}             % mudanca recente e vigilancia
\newcommand{\iAgir}{\faAmbulance}            % acao imediata e disposicao
\newcommand{\iProteger}{\faHandsHelping}     % suporte, protecao, continuidade
\newcommand{\iNorma}{\faBalanceScale}        % normativa, documento, registro
\newcommand{\iFarmaco}{\faPills}             % farmacologia

% As seis caixas, cada uma ja com o seu icone no titulo.
\newenvironment{caixaavaliar}[1]{\begin{CRbaseavaliar}{\iAvaliar~~#1}}{\end{CRbaseavaliar}}
\newenvironment{caixamudanca}[1]{\begin{CRbasemudanca}{\iMudanca~~#1}}{\end{CRbasemudanca}}
\newenvironment{caixaagir}[1]{\begin{CRbaseagir}{\iAgir~~#1}}{\end{CRbaseagir}}
\newenvironment{caixaproteger}[1]{\begin{CRbaseproteger}{\iProteger~~#1}}{\end{CRbaseproteger}}
\newenvironment{caixanorma}[1]{\begin{CRbasenorma}{\iNorma~~#1}}{\end{CRbasenorma}}
\newenvironment{caixafarmaco}[1]{\begin{CRbasefarmaco}{\iFarmaco~~#1}}{\end{CRbasefarmaco}}

% Titulo de frame com icone de orientacao. O icone marca a etapa da cadeia
% avaliar, formular, decidir, construir seguranca, garantir continuidade, e
% permite ao aluno saber onde esta sem ler o titulo inteiro.
\newcommand{\ftit}[2]{#1\,\,#2}

% =============================================================================
% REALCES DE TEXTO
% =============================================================================
\newcommand{\avaliar}[1]{\textbf{\color{campoazul!90!black}#1}}
\newcommand{\mudanca}[1]{\textbf{\color{campoambar!90!black}#1}}
\newcommand{\agir}[1]{\textbf{\color{campobordo!90!black}#1}}
\newcommand{\proteger}[1]{\textbf{\color{campoverde!80!black}#1}}
\newcommand{\norma}[1]{\textbf{\color{campografite}#1}}
\newcommand{\farmaco}[1]{\textbf{\color{camporoxo!88!black}#1}}

% Fala literal a ser dita ao paciente ou a familia. Vai para a tela como esta.
\newcommand{\fala}[1]{{\itshape\color{campoazul!85!black}\guillemotleft #1\guillemotright}}

% Procedencia discreta, ao pe do frame.
\newcommand{\fonte}[1]{%
  \par\vspace{0.4ex}%
  {\raggedleft\scriptsize\sffamily\color{campografite}#1\par}%
}

% Rotulo curto embutido no texto.
\newcommand{\chip}[2]{%
  \tikz[baseline=(x.base)]{\node[chipno={#1}] (x) {\strut #2};}%
}

% Grau de vinculacao da afirmacao. A aula separa o que e norma, o que e
% interpretacao clinica e o que e recomendacao apoiada em evidencia.
\newcommand{\selonorma}{{\scriptsize\sffamily\color{campografite}\textbf{[norma]}}}
\newcommand{\seloevidencia}{{\scriptsize\sffamily\color{campoverde}\textbf{[evid\^encia]}}}
\newcommand{\seloclinico}{{\scriptsize\sffamily\color{campoazul}\textbf{[interpreta\c{c}\~ao cl\'inica]}}}

% =============================================================================
% MENSAGEM DOMINANTE DO FRAME
% Largura total. Nunca dentro de columns.
% =============================================================================
\newcommand{\mensagem}[1]{%
  \par\vspace{\CRgapbox}%
  \begingroup
  \setlength{\fboxsep}{5pt}%
  \noindent\colorbox{campoazul!10}{%
    \parbox{\dimexpr\textwidth-2\fboxsep\relax}{%
      \centering\small\sffamily\bfseries\color{campoazul!90!black}#1}}%
  \endgroup
  \par
}
\newcommand{\mensagemagir}[1]{%
  \par\vspace{\CRgapbox}%
  \begingroup
  \setlength{\fboxsep}{5pt}%
  \noindent\colorbox{campobordo!8}{%
    \parbox{\dimexpr\textwidth-2\fboxsep\relax}{%
      \centering\small\sffamily\bfseries\color{campobordo}#1}}%
  \endgroup
  \par
}

% =============================================================================
% ITENS COMPACTOS, com a separacao governada pelo botao de ar
% =============================================================================
\newenvironment{itemc}{%
  \begin{itemize}\setlength{\itemsep}{\CRitemsep}\setlength{\parskip}{0pt}%
  \setlength{\topsep}{0pt}\setlength{\partopsep}{0pt}%
}{%
  \end{itemize}%
}

% Calha explicita entre colunas. Nenhum par de colunas soma mais de 0.94.
\newcommand{\calha}{\hspace{0.035\textwidth}}

% =============================================================================
% TABELAS DA SERIE
% -----------------------------------------------------------------------------
% A altura de linha vem de uma regua invisivel IDENTICA no inicio de TODA
% celula, governada por \CRestica, e nunca de \extrarowheight, que acrescenta
% espaco so acima da linha de base e desalinha o rotulo da primeira coluna.
% A primeira coluna e coluna de paragrafo de largura fixa, tipo E, e nunca l.
% Dentro do ambiente vale \lefthyphenmin=3. A barra de celula se escreve
% \slash, que imprime a mesma barra e admite quebra depois dela.
% =============================================================================
\newcommand{\CRregua}{\rule[-\CRestica]{0pt}{\dimexpr 2\CRestica + 1em\relax}}
\newcolumntype{E}{>{\CRregua\raggedright\arraybackslash}p{\CRcolentidade}}
\newcolumntype{P}[1]{>{\CRregua\raggedright\arraybackslash}p{#1}}
\newcolumntype{C}[1]{>{\CRregua\centering\arraybackslash}p{#1}}

\newenvironment{CRtabela}{%
  \begingroup
  \lefthyphenmin=3\relax
  \sffamily\footnotesize
  \setlength{\tabcolsep}{3pt}%
  \renewcommand{\arraystretch}{1.0}%
  \arrayrulecolor{campoazul!35}%
}{%
  \endgroup
}

% =============================================================================
% \quadroDx : macro UNICA dos frames da parte 3B. Geometria identica em todas
% as ocorrencias, dois quadros por frame. Nao revisa o diagnostico.
%   #1 quadro | #2 o que eleva o risco aqui | #3 o que muda na conduta
%   #4 erro classico
% =============================================================================
\newcommand{\quadroDx}[4]{%
  \begin{tcolorbox}[
      enhanced, sharp corners=downhill, arc=3pt, boxrule=0.9pt,
      colframe=campoazul!75!black, colback=white,
      left=5pt, right=5pt, top=4pt, bottom=4pt,
      width=\linewidth, height=6.15cm, valign=top,
      drop shadow={black!35},
      fontupper=\scriptsize,
      title={\sffamily\bfseries #1}, fonttitle=\scriptsize,
      colbacktitle=campoazul!88, coltitle=white,
      attach boxed title to top left={xshift=5pt, yshift=-2pt},
      boxed title style={sharp corners=downhill, arc=2pt, boxrule=0pt}]
    \vspace{0.26cm}
    {\sffamily\bfseries\scriptsize\color{campoambar!88!black}O QUE ELEVA O RISCO AQUI}\par
    \vspace{0.05cm}#2\par
    \vspace{0.17cm}
    {\sffamily\bfseries\scriptsize\color{campobordo}O QUE MUDA NA CONDUTA}\par
    \vspace{0.05cm}#3\par
    \vspace{0.17cm}
    {\sffamily\bfseries\scriptsize\color{campografite}ERRO CL\'ASSICO}\par
    \vspace{0.05cm}#4
  \end{tcolorbox}%
}
